
power. For example, on the topos @@M0@@ ringed by the field @@M1@@, the representation @@M2@@ defined by
@@M3@@
is not a power. In particular, 6.2.3 does not remain valid for an arbitrary locally ringed topos @@M4@@ and an arbitrary representation @@M5@@.

\sgaStableTarget{sga6.e10.para.6-5-3}\textbf{6.5.3.} On the other hand, 6.2.3 remains valid for the representations associated with formulas (6.2.8), (6.2.9) and (6.2.10) on an arbitrary ringed topos @@M6@@, because these representations are ``algebraic'' in the following sense:

Let @@M7@@ be the final object of @@M8@@, put @@M9@@, and consider the functor
@@M10@@
which assigns to each @@M11@@ the affine scheme @@M12@@ over @@M13@@. Then, for every positive integer @@M14@@, the presheaf @@M15@@ on @@M16@@ is the composite with @@M17@@ of the corresponding presheaf @@M18@@ on @@M19@@. A representation on @@M20@@
\sgaStableTarget{sga6.e10.eq.6-5-4}@@M21@@
is called \emph{algebraic} if it arises via @@M22@@ from a representation
\sgaStableTarget{sga6.e10.eq.6-5-5}@@M23@@
of presheaves of groups on @@M24@@, i.e., if one has
@@M25@@
for every @@M26@@, @@M27@@.

The representation @@M28@@ in (6.5.4) is, of course, completely determined by @@M29@@. It follows immediately that (6.2.3) remains valid in an arbitrary ringed topos when @@M30@@ is algebraic. In particular, formulas (6.2.8), (6.2.9) and (6.2.10) therefore remain valid on an arbitrary ringed topos.
